\color{uniblue}\section[Описание и работа]{Описание и работа\color{white!0}<ABBRS>} % вся эта белиберда только чтоб невидимый тэг воткнуть!
\color{black}

\color{uniblue}{\subsection{Назначение}}
\color{black}

\begin{enumerate}[label=\arabic{section}.\arabic{subsection}.\arabic*, labelsep=4pt, leftmargin=0pt, itemindent=57pt]

\item
Устройство серии <<ЮНИТ-М300-Т>> разработано для применения в схемах вторичной коммутации распределительных устройств 6-35 кВ с постоянным, переменным и выпрямленным оперативным током. Устройство предназначено для реализации функций релейной защиты силового трансформатора с высшим напряжением 35 кВ, при этом может осуществлять функции измерения, регистрации, осциллографирования и сигнализации. При подключении устройства к информационной сети энергообъекта возможно осуществление функций дистанционного управления устройством и сбор данных с него посредством протоколов и интерфейсов, определяемыми конфигурацией устройства. 
\item
Устройство предлагается с типовой конфигурацией: в технической части определяемой картой заказа (см. приложение \ref{app:card}), в функциональной части -- функционально-структурной схемой (см. приложение \ref{app:fsu}). Схемы подключения устройства по аналоговым входам приведены в приложении \ref{app:analogs}, схемы типового подключения сигналов алгоритмов к дискретным входам и релейным выходам устройства приведены в приложениях \ref{app:inouts}, \ref{app:inouts2}.
\item
По желанию Заказчика возможно выполнение устройства в других конструктивных исполнениях в соответствии с общей картой заказа на серию <<ЮНИТ-М300>>, которая приводится в документации ЮТКБ.656122.609 РЭ <<Устройство защиты и автоматики серии <<ЮНИТ-М300>>. Руководство по эксплуатации общее>>. 
\item
Устройство серии <<ЮНИТ-М300-Т>> разработано для применения в составе шкафов типа ШЭТ 111.12-1, ШЭТ 111.02-0 для установки на объектах филиалов и ДО ПАО <<Россети>>, построенных на базе архитектур I и II типа, а также на базе смешанной архитектуры.

\end{enumerate}

%========================== Технические характеристики устройства =====================

\color{uniblue}\subsection{Технические характеристики}
\color{black}

\begin{enumerate}[label=\arabic{section}.\arabic{subsection}.\arabic*, labelsep=4pt, leftmargin=0pt, itemindent=57pt]

\item
Основные номинальные параметры устройства указаны в таблице \ref{desc:chars}. 

\setlength{\extrarowheight}{0.06cm} % высота ячеек в таблице
\small
\begin{longtable}
    {|>{\centering\arraybackslash}m{12.9cm}|>{\centering\arraybackslash}m{4cm}|}
    \caption{Основные номинальные параметры устройства\hfill\vspace{-0.5\baselineskip}}\label{desc:chars}\\ 
    \hhline{|-|-|}
    \rowcolor{gray!30}
    Параметр & Значение \\ 
    \hhline{|-|-|}
    \endfirsthead

    %\caption{\emph{Продолжение\hfill\vspace{-0.5\baselineskip}}} \\ % Отступ обозначения таблицы от таблицы и выравнивание надписи слева
    \caption*{\hspace{3pt}\emph{Продолжение таблицы \ref{desc:chars}\hfill\vspace{-0.5\baselineskip}}} \\ % сделано по ГОСТ 2.105 п.6.8.7
    \hhline{|-|-|}
    \rowcolor{gray!30}
    Параметр & Значение \\ 
    \hhline{|-|-|}
    \endhead
    \hhline{|-|-|}

    \endfoot
    \endlastfoot

    \raggedright Номинальный переменный ток, А  & \raggedright\arraybackslash 1; 5 \\  \hhline{|-|-|}
    \raggedright Номинальное переменное фазное напряжение, В  & \raggedright\arraybackslash 100/$\sqrt{3}$ \\  \hhline{|-|-|}   
    \raggedright Номинальная частота, Гц  & \raggedright\arraybackslash 50 \\  \hhline{|-|-|}
    \raggedright Номинальное напряжение оперативного постоянного тока, В  & \raggedright\arraybackslash 110; 220 \\  \hhline{|-|-|}
    \raggedright Номинальное напряжение оперативного переменного тока, В  & \raggedright\arraybackslash 110; 220 \\  \hhline{|-|-|}
    %\flushleft \vspace{3pt} Период от подачи питания оперативного тока на МП устройство до начала выполнения им основных функций (время готовности устройства РЗА к срабатыванию), не более, с  & \flushleft\arraybackslash 1 \\  \hhline{|-|-|}

\end{longtable}
\normalsize
\setlength{\extrarowheight}{0.0cm} % возврат параметра высота ячеек в таблице

\end{enumerate}
%================================================================ Состав и конструктивное исполнение ==================================================

\color{uniblue}\subsection{Состав и конструктивное исполнение}
\color{black}

\begin{enumerate}[label=\arabic{section}.\arabic{subsection}.\arabic*, labelsep=4pt, leftmargin=0pt, itemindent=57pt]
    
\item
Корпус устройства в типовой конфигурации предусматривает одно исполнение по ширине, габаритные размеры устройства приведены в приложении \ref{app:sizes}. Внешний вид устройства для различных типовых исполнений показан на рисунках \ref{device:img1}, \ref{device:img2}. Подробное описание конструкции и функций используемых модулей (а также доступных к заказу) приведено в документации ЮТКБ.656122.609 РЭ <<Устройство защиты и автоматики серии <<ЮНИТ-М300>>. Руководство по эксплуатации общее>>.

\vspace{3mm}
\begin{figure}[H]
\centering
\includegraphics[width=1\textwidth,height=1\textheight,keepaspectratio]{Т_1арх.jpg}
\captionof{figure}{Вид устройства <<ЮНИТ-М319-Т>> с размещением модулей по слотам (исполнение для I архитектуры)}\label{device:img1}
\end{figure}

\vspace{3mm}
\begin{figure}[H]
\centering
\includegraphics[width=1\textwidth,height=1\textheight,keepaspectratio]{Т_2арх.jpg}
\captionof{figure}{Вид устройства <<ЮНИТ-М319-Т>> с размещением модулей по слотам (исполнение для II архитектуры)}\label{device:img2}
\end{figure}


%\color{uniblue}{\subsection{Функциональный состав}}
%\color{black}

\item
Перечень функций, реализованных в составе устройства, приведен в таблице \ref{desc:tbl1}. Описание основных функций устройства приводится в разделе \ref{sec:funcs} данной документации. Подробное описание вспомогательных (сервисных) функций приведено в документации ЮТКБ.656122.609 РЭ <<Устройство защиты и автоматики серии <<ЮНИТ-М300>>. Руководство по эксплуатации общее>>.

\setlength{\extrarowheight}{0.05cm} % высота ячеек в таблице
\small
\begin{longtable}
    {|>{\centering\arraybackslash}m{3.2cm}|>{\centering\arraybackslash}m{13.7cm}|}
    \caption{Функции в составе устройства <<ЮНИТ-М300-Т>>\hfill\vspace{-0.5\baselineskip}}\label{desc:tbl1}\\ 
    \hhline{|-|-|}
    \rowcolor{gray!30}
    Группа функций & Примечание \\ 
    \hhline{|-|-|}
    \endfirsthead

    \caption*{\hspace{3pt}\emph{Продолжение таблицы \ref{desc:tbl1}\hfill\vspace{-0.5\baselineskip}}} \\ % сделано по ГОСТ 2.105 п.6.8.7
    \hhline{|-|-|}
    \rowcolor{gray!30}
    Группа функций & Примечание \\ 
    \hhline{|-|-|}
    \endhead
    \hhline{|-|-|}

    \endfoot
    \endlastfoot

    \centering Основные функции & \flushleft\arraybackslash Газовые защиты трансформатора, технологические защиты трансформатора, резервные защиты трансформатора (МТЗ), сигнализация и управление выключателем стороны ВН трансформатора, УРОВ ВН  \\  \hline

    \multirow{7}{*}{\centering Сервисные функции} & \flushleft\arraybackslash Сигнализация (обеспечение ввода служебных дискретных сигналов о состоянии оборудования шкафа РЗА, формирование отчетов по изменению состояния дискретных внешних и внутренних сигналов и выдачу отчетов в АСУ~ТП, формирование и выдачу сигналов аварийной и предупредительной сигнализации в схему резервной сигнализации) \\ \cline{2-2}
    & \flushleft\arraybackslash  Осциллографирование и журналирование (запись доаварийных, аварийных и послеаварийных режимов с сохранением осциллограмм в энергонезависимой памяти устройства, запись о событиях по изменению конфигурации устройства, изменения состояния внутренних сигналов алгоритмов, обновление встраиваемого программного обеспечения и пр. с метками времени и описанием события в энергонезависимой памяти устройства) \\ \cline{2-2}
    & \flushleft\arraybackslash Самодиагностика (непрерывный оперативный контроль работоспособности в течение всего времени работ) \\ \cline{2-2}
    & \flushleft\arraybackslash  Конфигурирование (задание/переключение уставок, ввод в работу / вывод из работы, изменение параметров самоописания устройства, перезагрузка устройства, ввод/вывод тестового режима, изменение режима работы информационных сервисов, возможность переопределения связей входных и выходных логических сигналов алгоритмов устройства с дискретными входами, выходными реле, светодиодными индикаторами и функциональными клавишами управления) \\ \cline{2-2}
    & \flushleft\arraybackslash  Тестовый режим (обработка принятых сигналов и отправка отчетов с установленным флагом <<Тест>>) \\ \cline{2-2}
    & \flushleft\arraybackslash  Самоописание (предоставление данных об устройстве посредством устройства <<ЮНИТ-ИЧМ>> и сервисную программное обеспечение <<ЮНИТ-СЕРВИС>>) \\ \cline{2-2}
    & \flushleft\arraybackslash  Синхронизация времени (календарь и часы реального времени, синхронизируемые от персонального компьютера или от АСУ~ТП) \\ \hline

\end{longtable}
\normalsize
\setlength{\extrarowheight}{0.0cm} % возврат параметра высота ячеек в таблице

\end{enumerate}

%============================================================ Устройство и работа ==================================================
\color{uniblue}\subsection{Устройство и работа}\color{black}

\begin{enumerate}[label=\arabic{section}.\arabic{subsection}.\arabic*, labelsep=4pt, leftmargin=0pt, itemindent=57pt]
\item
Описание функций релейной защиты и автоматики в составе устройства, а также их функциональные схемы и уставки приводятся в разделе \ref{sec:funcs}. Общая функционально-структурная схема взаимосвязей функций, дискретных входов и выходов устройства приведена в приложении \ref{app:fsu}. Принятые обозначения элементов схем приведены в приложении \ref{app:func_desc}.
\item
Система самодиагностики устройства выполняет тесты в полном объеме при включении устройства, а также непрерывно в фоновом режиме.
Системой самодиагностики устройства в автоматическом режиме контролируются следующие параметры:
\begin{itemize}
\item состояние аппаратной части устройства, в том числе: АЦП модулей ввода аналоговых сигналов, блока питания, ОЗУ, ПЗУ, процессора, модулей дискретных входов, модулей дискретного управления, вспомогательных модулей;
\item температурный режим устройства;
\item состояние синхронизации времени;
\item сохранность исполняемого программного кода (целостность ПО);
\item состояние измерительных цепей;
\item исправность используемых каналов связи.
\end{itemize}
\item
Система самодиагностики фиксирует критические неисправности (<<Неисправность>>) и некритические ошибки (<<Ошибка>>). В случае появления критической неисправности активируется внутренний сигнал <<Критическая неисправность>>, замыкаются контакты реле IRF, прекращается работа алгоритмов основных функций в составе устройства, это состояние устройства сигнализируется свечением светодиода <<Неисправность>> и отсутствием свечения светодиода <<Готовность>>. В случае появления некритической ошибки активируется внутренний сигнал <<Некритическая ошибка>>, алгоритмы основных функций в составе устройства продолжают работать, это состояние устройства сигнализируется свечением светодиода <<Неисправность>> и отсутствием свечения светодиода <<Готовность>>. Подробное описание системы самодиагностики приводится в документации ЮТКБ.656122.609 РЭ <<Устройство защиты и автоматики серии <<ЮНИТ-М300>>. Руководство по эксплуатации общее>>.

\item
Для типоисполнения <<ЮНИТ-М314>> предусматривается установка модулей в любой из слотов с <<М3>> по <<М7>>:
\begin{itemize}
    \item выходных реле типа <<К001>> или <<К002>>;
    \item дискретных входов типа <<В001>> или <<В021>>.
\end{itemize}
\item
Для типоисполнения «ЮНИТ-М319» предусматривается установка модулей в любой из слотов с <<М3>> по <<М11>>:
\begin{itemize}
    \item выходных реле типа <<К001>>, <<К002>>;
    \item дискретных входов типа <<В001>> или <<В021>>.
\end{itemize}
\item
Типовое подключение входных сигналов функций к дискретным входам и выходным реле в составе устройства показано в приложениях \ref{app:inouts} и \ref{app:inouts2}. 

\item
В устройстве устанавливается измерительный модуль в зависимости от карты заказа, например модуль <<М044>> (в котором предусмотрено четыре канала для измерения переменного тока и четыре канала для измерения переменного напряжения) или модуль <<М030>> (в котором предусмотрено три канала для измерения переменного тока), в устройстве <<ЮНИТ-М314>> занимает место в слотах <<М08>>, <<М09>>, а в устройстве <<ЮНИТ-М319>> занимает место в слотах <<М12>>, <<М13>>. Типовое подключение входных аналоговых сигналов к модулю аналоговых входов в составе устройства показано в приложении \ref{app:analogs}.

\item
Подробное описание параметров настроек и технических характеристик модулей дискретных входов, модулей дискретного управления и измерительных модулей приводится в документации ЮТКБ.656122.609 РЭ <<Устройство защиты и автоматики серии <<ЮНИТ-М300>>. Руководство по эксплуатации общее>>.

\item
В устройствах с поддержкой протокола передачи данных МЭК 61850 в качестве входных сигналов ФСУ могут быть использованы данные из GOOSE-сообщений, на которые оформлена подписка. Количество сигналов GOOSE, получаемых от интеллектуальных устройств, в данном типоисполнении устройства равно 21 шт.
\item
Сигналы из GOOSE-сообщений проходят предварительную обработку перед тем как попасть в <<Блок объединительный>> и далее в ФСУ (см. приложение \ref{app:fsu}).
В предварительную обработку входят следующие операции:
\begin{itemize}
\item проверка качества принимаемого сигнала;
\item алгоритм обработки резервирования GOOSE-сообщений;
\item алгоритм обработки симуляции GOOSE-сообщений;
\item алгоритм обработки метки <<Test>> GOOSE-сообщений;
\item подстановка значения по умолчанию при <<плохом>> качестве данных, принимаемых из GOOSE-сообщений.
\end{itemize}
Подробнее о подписке на GOOSE-сообщения указано в RU.37182817.00001.02.34.01 – <<Программное обеспечение мониторинга и параметрирования <<ЮНИТ-СЕРВИС>>. Руководство пользователя>>.

\item 
В устройствах с поддержкой протокола передачи данных МЭК 61850 для связи со смежными устройствами может быть использован протокол GOOSE по МЭК 61850-8.1. Количество сигналов GOOSE-сообщений, передаваемых в смежные интеллектуальные устройства, зависит от типоисполнения ЮНИТ-М300, первичного оборудования и в общем случае составляет не менее 100 шт. Подробнее о публикации GOOSE-сообщений указано в RU.37182817.00001.02.34.01 – <<Программное обеспечение мониторинга и параметрирования <<ЮНИТ-СЕРВИС>>. Руководство пользователя>>.
\item
Устройство поддерживает два режима оперативного управления: местный и дистанционный. Выбор режима управления устройством производится посредством функциональной клавиши <<М/Д>> или по состоянию выбранного дискретного входа. В местном режиме игнорируются все команды управления с верхнего уровня АСУ ТП. Более подробно режимы устройства описаны в документации ЮТКБ.656122.609 РЭ <<Устройство защиты и автоматики серии <<ЮНИТ-М300>>. Руководство по эксплуатации общее>>.
\item
Для управления режимами работы основных функций (оперативный ввод/вывод, перевод на сигнал и пр.), а также для генерации различных сигналов сброса функций (сброс триггеров фиксации неисправности цепей газовых и технологических защит, сброс счетчиков и пр.), применяются виртуальные ключи и кнопки. Подробное описание и сценарии работы виртуальных ключей и кнопок приведено в документации ЮТКБ.656122.609 РЭ <<Устройство защиты и автоматики серии <<ЮНИТ-М300>>. Руководство по эксплуатации общее>>. 
\item
В данном типоисполнении устройства реализовано четыре группы уставок. Подробное описание работы с группами уставок приведено в документации ЮТКБ.656122.609 РЭ <<Устройство защиты и автоматики серии <<ЮНИТ-М300>>. Руководство по эксплуатации общее>>.
\item
Устройство предусматривает возможность синхронизации внутренних часов реального времени несколькими способами:
\begin{itemize}
    \item по протоколу SNTP;
    \item по протоколу PTP;
    \item через дискретный вход устройства с использованием сигнала синхронизации 1PPS;
    \item {средствами стандартизованного протокола передачи данных.}
\end{itemize}
Подробное описание способов синхронизации времени устройства приведено в документации ЮТКБ.656122.609 РЭ <<Устройство защиты и автоматики серии <<ЮНИТ-М300>>. Руководство по эксплуатации общее>>.
\item
Устройство предусматривает возможность ведения журнала событий. Перечень сигналов, доступных для регистрации в журнале событий, приведен в приложении \ref{app:signals}. События безопасности записываются в отдельный журнал безопасности. Подробное описание параметров журналов событий приведено в документации ЮТКБ.656122.609 РЭ <<Устройство защиты и автоматики серии <<ЮНИТ-М300>>. Руководство по эксплуатации общее>>.
\item
Устройство предусматривает возможность осциллографирования событий. Перечень дискретных сигналов, доступных для регистрации аварийным осциллографом, приведен в приложении \ref{app:signals}. Подробное описание параметров встроенного регистратора аварийных событий приведено в документации ЮТКБ.656122.609 РЭ <<Устройство защиты и автоматики серии <<ЮНИТ-М300>>. Руководство по эксплуатации общее>>.
\item Перечень аналоговых сигналов, доступных для регистрации аварийным осциллографом, приведен в таблице \ref{osz:tbl1}.

\setlength{\extrarowheight}{0.05cm}
\small
\begin{longtable}{|>{\centering\arraybackslash}m{6.8cm}|>{\centering\arraybackslash}m{2cm}|>{\centering\arraybackslash}m{5.3cm}|>{\centering\arraybackslash}m{2cm}|}
\caption{Перечень аналоговых сигналов, доступных для регистрации аварийным осциллографом\hfill\vspace{-0.5\baselineskip}}\label{osz:tbl1}\\ 
\hline
\rowcolor{gray!30}
Описание аналогового сигнала & Канал & Наименование  & Пуск РАС \\ 
\hline
\endfirsthead
\caption*{\hspace{3pt}\emph{Продолжение таблицы \ref{osz:tbl1}\hfill\vspace{-0.5\baselineskip}}} \\ % сделано по ГОСТ 2.105 п.6.8.7
\hline
\rowcolor{gray!30}
Описание аналогового сигнала & Канал & Наименование  & Пуск РАС \\ 
%\hline
\endhead
%\hline
\endfoot
%\hline
\endlastfoot

\raggedright Фазные токи ВН & \centering 1 & \centering Ia\_{\text{ВН}}  & \centering\arraybackslash + \\
\hline
\raggedright >>             & \centering 2 & \centering Ib\_{\text{ВН}}  & \centering\arraybackslash + \\
\hline
\raggedright >>             & \centering 3 & \centering Ic\_{\text{ВН}}  & \centering\arraybackslash + \\
\hline

\raggedright Вычисленный ток ОП ВН & \centering 4 & \centering I2\_{\text{ВН}}  & \centering\arraybackslash + \\
\hline
\raggedright Вычисленный ток НП ВН & \centering 5 & \centering 3I0\_{\text{ВН}}  & \centering\arraybackslash + \\
\hline

\raggedright Фазные напряжения НН & \centering 6 & \centering Ua\_{\text{НН}$^{*)}$}  & \centering\arraybackslash - \\
\hline
\raggedright >>                    & \centering 7 & \centering Ub\_{\text{НН}$^{*)}$}  & \centering\arraybackslash - \\
\hline
\raggedright >>                    & \centering 8 & \centering Uc\_{\text{НН}$^{*)}$}  & \centering\arraybackslash - \\
\hline

\raggedright Вычисленное напряжение ОП НН & \centering 9 & \centering U2\_{\text{НН}$^{*)}$}  & \centering\arraybackslash - \\
\hline
\raggedright Вычисленное напряжение НП НН & \centering 10 & \centering 3U0\_{\text{НН}$^{*)}$}  & \centering\arraybackslash - \\
\hline

% Добавление строки внизу таблицы
\multicolumn{4}{|l|}{$^{*)}$ Только для исполнения устройства, предназначенного для I архитектуры.} \\
\hline


\end{longtable}
\normalsize
\setlength{\extrarowheight}{0.0cm}

\item В устройстве реализована поддержка режима тестирования <<Тест>> в объеме требований стандартов МЭК 61850. Особенности перевода устройства в режим тестирования и функциональности устройства в этом режиме описаны в документации ЮТКБ.656122.609 РЭ <<Устройство защиты и автоматики серии <<ЮНИТ-М300>>. Руководство по эксплуатации общее>>.

\item\sloppy Описание и настройки доступных портов и интерфейсов связи приводится в документации ЮТКБ.656122.609 РЭ <<Устройство защиты и автоматики серии <<ЮНИТ-М300>>. Руководство по эксплуатации общее>>.

\item Устройство имеет встроенную реализацию правил разграничения прав доступа, для этого в устройстве существует ролевая модель разграничения прав доступа.
В устройстве предусмотрены следующие роли:
\begin{itemize}
    \item Гость (пароль не требуется, доступен просмотр данных устройства);
    \item Дежурный;
    \item Инженер;
    \item Офицер (сотрудник) безопасности;
    \item Администратор (пароль <<АААА>> (кириллица), учетная запись предназначена для первоначального конфигурирования штатного профиля пользователей).
\end{itemize}
Все пользователи делятся на непривилегированных (роль Гость) и привилегированных (все остальные роли).
\item
Предустановленные пользователи с именами Гость и Администратор не подлежат удалению и принудительной блокировке, их роли не подлежат изменению. Гость используется как роль по умолчанию для доступа к просмотру данных устройства в режиме чтение без запроса аутентификационных данных пользователя. Администратор необходим для добавления/удаления других пользователей, установки им соответствующих ролей и установке/изменении паролей пользователей.
\item
После перезагрузки устройства подключенная сессия привилегированного пользователя сбрасывается. Устройство запускается с сеансом непривилегированного пользователя (роль Гость).
Подробное описание ролей и доступных полномочий и ограничений для каждой из ролей приведено в документации ЮТКБ.656122.609 РЭ <<Устройство защиты и автоматики серии <<ЮНИТ-М300>>. Руководство по эксплуатации общее>>.

\end{enumerate}

%============================================================ Средства измерения, инструмент и принадлежности ==================================================
\color{uniblue}\subsection{Средства измерения, инструмент и принадлежности}\color{black}
Перечень оборудования и средств измерения, рекомендуемый для проведения эксплуатационных проверок устройства, приведен в документации ЮТКБ.656122.609 РЭ <<Устройство защиты и автоматики серии <<ЮНИТ-М300>>. Руководство по эксплуатации общее>>.

%================================================================= Маркировка и пломбирование ==================================================
\color{uniblue}\subsection{Маркировка и пломбирование}\color{black}
Сведения по маркировке и пломбировании устройства приведены в документации ЮТКБ.656122.609 РЭ <<Устройство защиты и автоматики серии <<ЮНИТ-М300>>. Руководство по эксплуатации общее>>.

%======================================================================== Упаковка ==================================================
\color{uniblue}\subsection{Упаковка}\color{black}
Сведения по упаковке устройства приведены в документации ЮТКБ.656122.609 РЭ <<Устройство защиты и автоматики серии <<ЮНИТ-М300>>. Руководство по эксплуатации общее>>.